\documentclass[11pt]{article}
\usepackage[margin=1in]{geometry}
\usepackage{graphicx}
\usepackage{booktabs}
\usepackage{amsmath}
\usepackage{microtype}
\usepackage{xurl}
\setlength{\emergencystretch}{3em}
\usepackage{caption}
\usepackage{float}
\graphicspath{{figures/}{../re_extension_section/figures/}}

\title{NIMBY and Housing Extensions: RE and Coalitions}
\author{}
\date{}

\begin{document}
\maketitle

\section{Extension 1: rational expectations}

The first extension asks whether the paper's demographic transition changes materially once
households are assumed to form model-consistent expectations over the future path of house prices.
The published benchmark keeps the political mechanism in place but assumes that households do not
forecast future changes in political support and instead treat future prices using the original
forecasting rule. The rational-expectations exercise therefore asks whether this shortcut is doing
substantive economic work.

The current implementation keeps the exercise deliberately narrow. It shuts down coalition politics
and solves a no-politics transition problem from 2010 to 2018 under a guessed future price path.
Given that path, households solve backward, the cross-sectional distribution is simulated forward,
and the implied market-clearing future price path is recovered. A rational-expectations equilibrium
requires the guessed path to coincide with the implied path. This turns the problem into a
high-dimensional backward-forward fixed point over an entire sequence of prices.

\begin{figure}[H]
\centering
\includegraphics[width=0.98\textwidth]{re_lambda_continuation.pdf}
\caption{Rational-expectations continuation across demographic shock sizes. The left panel plots
the final rational-expectations price path for $\lambda \in \{0,0.5,1\}$ as log deviations from
the baseline price level. The right panel reports the residual norm and maximum pointwise pricing
gap along the same continuation ladder. The key pattern is smoothness rather than a regime break.}
\end{figure}

Figure 1 summarizes the continuation exercise. As the full demographic transition is turned on, the
full-shock rational-expectations solution reconnects smoothly to nearby easier problems. This is
useful because it suggests that the extension is not hiding a qualitatively different
rational-expectations branch. Instead, the same localized numerical bottleneck reappears as the
full shock is restored. The hard region migrates back toward the third model period, but the path
itself remains close to the benchmark control-like solution.

\begin{figure}[H]
\centering
\includegraphics[width=0.82\textwidth]{re_vs_paper_benchmark.pdf}
\caption{Heuristic comparison between the published benchmark and the rational-expectations
extension over 2010--2018. The black line is digitized from the saved local benchmark figure in the
original project workspace and normalized at 2010. The orange dashed line is the full-shock
rational-expectations no-politics transition, also normalized at 2010. The comparison is
directional rather than exact because the benchmark retains politics while the extension shuts it
down.}
\end{figure}

Figure 2 is the practical comparison for the paper. Over the shared 2010--2018 window, the
published benchmark rises materially while the current no-politics rational-expectations path is
nearly flat. The present rational-expectations extension therefore does not appear to ``bring
NIMBY forward.'' The more defensible interpretation is narrower: replacing the paper's expectations
shortcut with a model-consistent future price path does not overturn the benchmark transition. For
paper purposes, the extension is therefore best read as a robustness result. Rational expectations
matter at the margin, but they do not appear to generate a fundamentally different housing-market
transition under the current calibration.

The computational difficulty is nevertheless informative. The full rational-expectations problem is
hard not because the machine is slow, but because the equilibrium object is an entire future price
path and each candidate path requires a full backward solution plus forward simulation. The
selection diagnostics show a sharp local bottleneck around period 3, and the continuation exercise
shows that this bottleneck returns as the full demographic transition is restored. The extension is
therefore useful both as a robustness check and as a transparent explanation of why a fully solved
transition-path rational-expectations benchmark is numerically demanding.

\section{Extension 2: coalitions}

The second extension is more tightly connected to the paper's core mechanism. Instead of changing
the expectations rule, it changes the political aggregation rule. The goal is to ask whether the
paper's anti-supply mechanism becomes stronger when politically aligned owner groups receive more
effective political weight than under equal headcount voting.

This is intentionally a reduced-form political amplification exercise rather than a full coalition
formation game. The benchmark political model is nested exactly when all coalition weights are set
to zero. From that benchmark, the extension reweights the voting influence of homeowners, older
homeowners, leveraged owners, and large-housing-position owners one margin at a time, and then in a
combined case. Because only the political weighting object changes, the comparison is clean: the
extension asks whether stronger owner-side political influence shifts the steady-state political
equilibrium in the expected direction.

\begin{figure}[H]
\centering
\includegraphics[width=0.98\textwidth]{coalition_extension.pdf}
\caption{Reduced-form coalition-power extension. The left panel reports the equilibrium price shift
relative to the equal-weight benchmark for each one-margin case and the combined coalition case. The
right panel plots the benchmark and combined coalition objective curves across the same house-price
grid. Coalition weighting shifts the minimum toward a higher equilibrium price, with the strongest
movement coming from owner-based political overweighting.}
\end{figure}

Figure 3 delivers the main coalition message. The equal-weight benchmark price is about 1.8138. The
owner-only, old-owner-only, and combined coalition cases all raise the selected equilibrium price to
about 1.8276, which is roughly a 0.76 percent increase on this grid. By contrast, the
leveraged-owner-only and big-house-only cases move the weighted vote object without changing the
selected price at this resolution. The practical implication is that the coalition force that seems
to matter most is not ``owners'' in every possible sense, but specifically the groups most closely
aligned with the paper's baseline mechanism: homeowners, especially older homeowners.

This changes the paper more than the rational-expectations extension, but it changes it in a
reinforcing way rather than an overturning way. The coalition-power exercise does not reverse the
paper's lifecycle NIMBY mechanism. It sharpens it. A natural paper-level interpretation is that the
benchmark mechanism may understate the anti-supply force if politically organized owner groups carry
more influence than equal headcount would imply. That makes the extension worth including, but it
should still be described honestly as a reduced-form coalition-power exercise rather than a full
coalition game.

A true coalition model would require substantially more structure: explicit group formation,
bargaining over policy, and dynamic political negotiation over time. That is a much richer object
than what is needed for the current paper. For present purposes, the reduced-form extension is
useful precisely because it isolates the amplification channel cleanly and shows that stronger owner
political power pushes the model in the same direction as the published benchmark.

\section{Interpretation}

Taken together, the two extensions push in different ways. The rational-expectations extension is a
supporting robustness exercise: it suggests that the paper's lifecycle NIMBY mechanism is not
primarily an artifact of the expectations shortcut. The coalition extension is more substantive: it
shows that giving more effective political weight to owner groups strengthens the anti-supply
equilibrium in the expected direction. If only one extension is emphasized in the paper, the
coalition-power result is the one closer to the core economic mechanism. If both are included, the
cleanest summary is that rational expectations leave the paper's message broadly intact, while
political amplification through owner-side coalition power strengthens it.

\end{document}
